\section{Related work}
\label{sec:related_work}

\paragraph{3D reconstruction of static scenes.}
The field of 3D reconstruction of static scenes, such as multi-view geometry~\citep{Ding:2022:TMV,Zhang:2023:GMV,Bae:2022:MVD}, structure from motion (SfM)~\citep{Pan:2024:GSF,Lindenberger:2021:PPS,Wang:2024:VSV}, and simultaneous localization and mapping (SLAM)~\citep{Lipson:2024:DPV,Teed:2021:DRO,Zhu:2024:NSL}, has been also significantly advanced by learning-based approaches.
Canonical approaches~\citep{Agarwal:2011:BRI,Furukawa:2010:TIS,Schonberger:2016:SFM} rely on a well-established pipeline based on projective geometry, such as feature extraction, triangulation, bundle adjustment, \etc.
However, learning-based approaches demonstrate that each component in the pipeline can be advanced by data-driven approaches.
Even the whole pipeline can be end-to-end~\citep{Weinzaepfel:2023:CRO,Wang:2024:DUS,Wang:2025:vgg,Zhang:2025:FLA,Wang:2024:3DR,Liang:2024:FFB,Elflein:2025:LTF} when networks learn strong priors from large-scale data.  
However, these methods specifically target static scenes that satisfy the epipolar geometry constraint.

\paragraph{3D reconstruction of dynamic scenes.}
Traditional approaches~\citep{Kumar:2017:MD3,Mustafa:2016:TC4,Luiten:2020:TTR} often tackle this problem using multi-stage pipelines that combine given depth and motion cues and optimize the main objective using rigidity or piecewise smoothness priors about geometry and motion.
Recent approaches demonstrate that this highly ill-posed problem can simply be solved by learning strong priors on geometry of dynamic scenes.
Those strong prior can be obtained by light-weight finetuning~\citep{Zhang:2025:MST,Lu:2024:A3R} of static 3D reconstruction model~\citep{Wang:2024:DUS}, large-scale training~\citep{Wang:2025:CUT}, or repurposing a generative model~\citep{Mai:2025:CVD}.
While these feedforward methods generate per-frame pointclouds, they lack temporal correspondence, limiting their application to motion understanding tasks.

\paragraph{Dense 4D reconstruction.}
Dense 4D reconstruction methods mainly fall into two categories.
One line of work relies on per-scene test-time optimization~\citep{Wang:2024:SM4,Wang:2024:TEE,Lei:2024:MOS,Liu:2025:MOD}. 
The approaches show high-fidelity reconstruction, but often take hours to run and depend on pre-computed cues such as camera poses and optical flow.
To enable real-time application, other feedforward approaches have been proposed~\citep{Han:2025:DDU,Sucar:2025:DPM,Liang:2025:ZSM,Feng:2025:ST4}, represented in sparse pointmaps.
Concurrent works are primarily designed for novel view synthesis and require camera poses at test time~\citep{lin2025movies,lin2025dgs,xu20254dgt}, while others require training separate heads for each downstream task~\citep{Badki:2025:L4P}.
In contrast, our method reconstructs a dense and explicit 4D representation in a feedforward manner without requiring camera poses as input.
Moreover, our holistic approach is suitable for both reconstruction and synthesis, demonstrating that synthesis improves reconstruction tasks. 

\paragraph{Dense 4D datasets.}
The lack of large-scale, densely annotated 4D datasets is a major bottleneck for the development of dense 4D reconstruction. 
While synthetic data~\citep{Zheng:2023:POA} suffers from a domain gap, the recent real-world Stereo4D dataset~\citep{Jin:2025:S4D} provides valuable supervision, though its annotations remain sparse and can be unreliable in occluded or distant regions.
\ourmethod{}, which uses a unified dynamic 3D Gaussian representation, effectively leverages self-supervision to overcome the limitations of sparse ground truth annotations in existing datasets.

